% Source-aligned English for printed p.196.
\clearpage
\SourcePage{196}

\LayerTarget{tohoku:result:abelian-g-sheaves-additive}
In particular, the abelian $G$-sheaves form an additive category, whose
properties we shall now indicate.
\LayerTarget{tohoku:definition:g-o-module}
More generally, let ${\mathbf O}$ be a $G$-sheaf of rings with identity, and consider
the sheaves $A$ on $X$ that are both abelian $G$-sheaves and ${\mathbf O}$-modules, the
operations of ${\mathbf O}$ on $A$ being compatible with the operations of $G$ (i.e.
the natural homomorphism of sheaves of sets
${\mathbf O}\times A\to A$ being a $G$-homomorphism). Such a sheaf will be
called a \textit{$G$-${\mathbf O}$-module}. Calling a sheaf homomorphism that is both an
${\mathbf O}$-homomorphism and a $G$-homomorphism a
\textit{$G$-${\mathbf O}$-homomorphism}, we see that the sum or composite of two
$G$-${\mathbf O}$-homomorphisms is a $G$-${\mathbf O}$-homomorphism.
\LayerTarget{tohoku:definition:category-g-o-modules}
Thus the $G$-${\mathbf O}$-modules form an \textit{additive category}, denoted
${\mathbf C}^{{\mathbf O}(G)}$; if ${\mathbf O}$ is the constant sheaf of integers (with the trivial
operations of $G$), we recover the category of abelian $G$-sheaves, denoted
${\mathbf C}^{X(G)}$.
\LayerTarget{tohoku:comparison:trivial-action-g-o-modules}
If $G$ acts trivially on $X$, the category ${\mathbf C}^{{\mathbf O}(G)}$ may be interpreted as
the category of ${\mathbf O}'$-modules, where ${\mathbf O}'$ is a suitable sheaf of rings. For
example, if $G$ acts trivially on ${\mathbf O}$, we have
${\mathbf O}'={\mathbf O}\otimes_{\mathbf Z}\mathbf Z(G)$, where $\mathbf Z(G)$ denotes the
algebra of $G$ over the ring $\mathbf Z$ of integers, or rather the constant
sheaf that it defines. The results of 3.1 remain valid with small
modifications:

\LayerTarget{tohoku:proposition:5.1.1}
\textsc{Proposition 5.1.1.} \textit{Let ${\mathbf O}$ be a $G$-sheaf of rings
on $X$. Then the additive category ${\mathbf C}^{{\mathbf O}(G)}$ of $G$-${\mathbf O}$-modules is an
abelian category satisfying axioms \textup{AB 5)} and \textup{AB 3*)} of \textup{1.5}, and it has a
generator.}\textsuperscript{9)}

\LayerTarget{tohoku:proof:5.1.1:p1}
The verification is immediate except for the existence of a generator, which
we shall now construct, generalizing the construction of 3.1.
\LayerTarget{tohoku:construction:generator-l-u}
For every open subset $U$ of $X$, let $L(U)$ be the ${\mathbf O}$-module that is the
direct sum of the ${\mathbf O}$-modules ${\mathbf O}_{g\cdot U}$ for $g\in G$. The sheaf $L(U)$
can evidently be regarded as a $G$-sheaf, so that it is even a
$G$-${\mathbf O}$-module. If $A$ is any $G$-${\mathbf O}$-module, a $G$-${\mathbf O}$-homomorphism from
$L(U)$ to $A$ is determined by its restriction to ${\mathbf O}_U$, which is an
${\mathbf O}$-homomorphism from ${\mathbf O}_U$ to $A$ and may, moreover, be prescribed
arbitrarily in advance. Giving such a homomorphism is also equivalent to
giving a section of $A$ over $U$.
\LayerTarget{tohoku:result:generator-family-g-o-modules}
If $B$ is a proper sub-$G$-${\mathbf O}$-module of $A$, then on some suitable open
subset $U$ one can find a section of $A$ that is not a section of $B$.
Consequently, the family of the $L(U)$ is a family of generators of
${\mathbf C}^{{\mathbf O}(G)}$, and their direct sum, as $U$ ranges over all open subsets of $X$,
is therefore a generator.

\LayerTarget{tohoku:conclusion:injective-g-o-embedding:p1}
It follows in particular from Proposition 5.1.1 that every $G$-${\mathbf O}$-sheaf is
isomorphic to a subsheaf of an injective $G$-${\mathbf O}$-module. It is important for
the

\SourceFootnote{9}{\LayerTarget{tohoku:footnote:9}
This proposition, as well as various subsequent results, can also be
stated in the following more general context. Let a category ${\mathbf C}$, a group
$G$, and a ``representation of $G$ by functors in ${\mathbf C}$'' be given: to every
$g\in G$ there is associated a functor $F_g:{\mathbf C}\to {\mathbf C}$, in such a
way that $F_e$ is the identity functor and $F_gF_{g'}=F_{gg'}$ (at least up to
\textit{specified} functorial isomorphisms satisfying certain coherence conditions
that we shall not make explicit). One can then construct the category ${\mathbf C}^G$
of ``objects of ${\mathbf C}$ with operator group $G$'', just as ${\mathbf C}^{X(G)}$ (for a
space with operator group $X(G)$) can be constructed from ${\mathbf C}^X$. (Compare
also the special case, Example 1.7.f). Many properties that hold for ${\mathbf C}$ are
then preserved upon passage to ${\mathbf C}^G$.}

\TohokuContinuation{197}
